% Generated draft: requires mathematical and source review.
\documentclass[11pt]{article}
\usepackage[margin=25mm]{geometry}
\usepackage{fontspec}
\setmainfont{Latin Modern Roman}
\usepackage{amsmath,amssymb,mathtools,unicode-math}
\setmathfont{Latin Modern Math}
\usepackage{enumitem,needspace,xurl,hyperref}
\setlength{\parindent}{0pt}
\setlength{\parskip}{5pt}
\setlength{\emergencystretch}{3em}
\newcommand{\N}{\mathbb{N}}
\newcommand{\Z}{\mathbb{Z}}
\newcommand{\Q}{\mathbb{Q}}
\newcommand{\R}{\mathbb{R}}
\newcommand{\C}{\mathbb{C}}
\newcommand{\F}{\mathbb{F}}
\newcommand{\A}{\mathbb{A}}
\newcommand{\PP}{\mathbb{P}}
\newcommand{\E}{\mathbb{E}}
\newcommand{\eps}{\varepsilon}
\newcommand{\dd}{\,\mathrm{d}}
\begin{document}
\section*{UnsolvedMath ID: 800472}
\textbf{Source problem code:} catalogue entry
\subsection{Definitions and mathematical statement}
The source poses a construction and theory-development problem rather than a conjecture with specified hypotheses. Here an infinite-dimensional billiard means an intended dynamical model with an infinite-dimensional configuration space, motion between collisions, and a rule governing reflection at a boundary. The source does not specify the configuration space, its geometric or analytic structure, the boundary, the admissible trajectories, or the reflection law; these are part of the requested construction.

Find physically interesting examples of infinite-dimensional billiards, giving precise definitions of their configuration spaces, free motion, collision configurations, and reflection dynamics. Possible connections with equations of mathematical physics are explicitly invited, but no particular equation or connection is required by the source. The phrase ``physically interesting'' is not given a mathematical criterion.

The motivation is a relationship between billiard and geodesic flows through appropriate limiting procedures. The source describes this relationship as folklore that becomes a theorem under certain conditions, without specifying those conditions or limits. It also recalls the Euler--Arnold interpretation of geodesics on infinite-dimensional groups. In particular, it describes solutions of the Euler equation of incompressible hydrodynamics as geodesics on the group of volume-preserving diffeomorphisms equipped with a right-invariant \(L^2\) metric. The underlying fluid domain, regularity, and boundary conditions are not specified. Similar frameworks are cited for the KdV equation, magnetohydrodynamics, and the Heisenberg chain, without further details.

An explicitly suggested, optional direction is to replace finitely many particles undergoing billiard reflection at a wall by a continuum of particles. The source points out that the reflected and incoming flows then occupy the same region, producing two distinct flows simultaneously through each point near the boundary in the proposed picture. It suggests using multiphase flows and a groupoid structure for generalized flows to accommodate this overlap. Develop, if possible, a theory of multiphase fluids with billiard reflections at the boundary; in the two-phase case, investigate whether such infinite-dimensional fluid billiards can be described by appropriate boundary conditions in a diffeomorphism groupoid. The relevant groupoid, boundary conditions, and precise meaning of reflection remain to be supplied. This suggested mechanism is a research direction, not an asserted construction or correspondence.
\subsection{Short English statement}
Construct mathematically defined infinite-dimensional billiards with physical significance, possibly connected to equations of mathematical physics. One suggested direction uses multiphase fluid flows and a diffeomorphism groupoid to describe overlapping incoming and reflected continua of particles. The source leaves the spaces, reflection laws, boundary conditions, and criterion of physical interest unspecified.
\subsection{Sources}
\begin{itemize}
\item Local UnsolvedMath registry, permanent ID \(800472\), source problem code \texttt{catalogue entry}; supplied complete source record from the AMR Open Problem Lists collection. The record is labelled \texttt{open}; this is supplied registry status, not an independently verified assessment.
\item Albach, Barraud et al., \emph{Open Problems in Billiards and Quantitative Symplectic Geometry} (2026), source item ``Boris Khesin --- Problem 1.'' Supplied source URL: \url{https://arxiv.org/abs/2602.12896}. The registry lists Albach, Barraud et al. as proposers and Boris Khesin as contributor. The supplied provenance records access on 2026-10-16 and extraction by \texttt{arxiv\_latex}; these are source-provided provenance claims, not a claim of independent access or verification here.
\item Source reference labels \texttt{[Giannetto, Lange]} are cited for relationships between billiard and geodesic flows through limits. Full bibliographic entries and the applicable hypotheses are absent from the supplied record.
\item Source reference labels \texttt{[Arnold66, ArnKh]} are cited for the infinite-dimensional geodesic and mathematical-physics background. Full bibliographic entries are absent from the supplied record.
\item Source reference labels \texttt{[Brenier, IzKh]} are cited for multiphase flows and the groupoid structure for generalized flows. Full bibliographic entries are absent from the supplied record.
\end{itemize}
\end{document}
